\section{Introduction}
Bolts, in particular, are among the most commonly used machine elements for component connection \cite{ste12} and are used wherever a detachable connection must reliably transmit load. Their mechanical function depends primarily on the preload generated during assembly. If the preload is insufficient, the clamped parts can separate earlier, the fraction of external load carried by the bolt increases, and the fatigue-critical stress amplitude in the bolt rises. Reliable preload assessment is therefore central to the design and maintenance of safety-relevant bolted joints. 
The nominal preload and load-bearing capacity of highly loaded bolted joints are commonly verified in accordance with VDI 2230 \cite{vdi2230}. The most critical joints typically bear the heaviest loads. The actual preload force, crucial for the bolted connection’s functionality, is achieved by applying the tightening torque during assembly. However, the tightening factor varies significantly across assembly methods, as determined by the assembly tool and the control measurement variable \cite{wie07}.

In practice, preload is often controlled indirectly by tightening torque. This is convenient but not a direct measure of bolt force, because a large portion of the applied torque is consumed by thread and bearing-surface friction\cite{vdi2230}. Variations in lubrication, coatings, surface roughness, repeated tightening, and assembly tool behavior can therefore lead to significant preload scatter. As a result, the actual preload state of an assembled joint may remain uncertain, even though it is the decisive functional quantity governing load transfer and fatigue behavior. Once the joint is assembled, the challenge becomes more severe: the preload is not directly visible, and direct measurement typically requires specialized hardware. This motivates inspection concepts that can estimate preload-related information from accessible, non-instrumented surfaces.

Established preload-measurement approaches include sensorized bolts, force-measuring washers, strain gauges, load cells, fiber-optic sensors, ultrasonic techniques, and bolt-length measurements. These methods are important for calibration, laboratory validation, and selected high-value applications, but they typically require additional hardware, wiring, coupling, calibration, or modification of the fastener or joint. For large numbers of standard fasteners, this limits scalability.

Computer vision has been investigated for bolt inspection, especially for detecting missing bolts, corrosion, visible looseness, or rotation of marked bolts \cite{afifi2024data}. These methods are useful for inspection, but most of them estimate a visible state such as bolt rotation or loosening angle rather than the mechanical preload itself \cite{cha2016vision,pham2020bolt,yu2021detection,sun2022yolo,qian2024loosening,malik2024review}. A bolt can lose a functionally relevant amount of preload before a large visible rotation or separation gap is detectable. This motivates a different visual target: not the rotation of the bolt, but the preload-induced surface deformation of the visible bolt-head face.

This paper proposes \emph{PreForceVision}, a physics-guided optical framework for estimating bolt preload from the deformation of the bolt-head front face. The central hypothesis is that axial preload produces a repeatable, micrometer-scale deformation and strain field on the bolt-head surface. If this field can be measured or approximated from reference-loaded image pairs, it can serve as a preload-sensitive optical signature. The proposed approach is a feasibility framework. It first asks whether the mechanical signal is observable, then how it can be extracted from images, and how it can be mapped to preload under controlled conditions.

%Preliminary results support the mechanical plausibility of the idea. For an M16 strength class 8.8 bolt under a preload of \SI{90}{\kilo\newton}, finite-element analysis indicates an approximately \SI{12}{\micro\meter} depression of the bolt-head face and an average surface deformation of about \SI{0.1}{\percent}. Strain-gauge measurements on a sensor bolt show comparable compressive strains, with an average value of approximately \SI{-0.116}{\percent} at \SI{90}{\kilo\newton}. A separate optical 3D comparison measurement using structured light shows a central depression of approximately \SIrange{10}{20}{\micro\meter} relative to the edge region after tightening. These observations indicate that the preload-induced deformation is physically measurable, while also showing that the expected optical signal is small and must be treated with strict measurement discipline.

\begin{figure*}[t]
   \centering
    \includegraphics[width=0.95\textwidth, keepaspectratio, trim=0.cm 0cm 0cm 0cm, clip]{Main/figures/fig1-framework.png}
    \caption{PreForceVision Framework: The proposed framework compares a reference and a loaded image of the bolt-head front face, extracts a preload-sensitive deformation signature, and maps this signature to preload with uncertainty. Hidden assembly geometry is not inferred directly from the head image.}
    \label{fig:concept}
\end{figure*}
